In this work, a bi-level optimization method is proposed for the solution of
multimodal design optimization problems with gradient-based local solvers. The
design space is divided into a Cartesian grid of boxes, and the selection of a
box is encoded by one categorical variable per design component. A mixed-integer
master problem selects the box to be explored from outer-approximation cuts
built on the boxes already solved, and the original problem is then solved
within the selected box by a local optimizer. The slopes of the cuts are
obtained by post-optimal sensitivity analysis of the sub-problem, and the size
of the master problem grows with the sum of the subdivisions instead of their
product. Since the value function is not convex in the box selection, the cuts
have to be corrected. A convexification constant and an adaptive repair of the
slopes are compared, and a sweep of the convexity value over the parallel probes
of the master is introduced to avoid its calibration. The influence of the
trust region of the master and of the subdivision density is also studied. The
method is assessed on four analytic test functions in two and five dimensions
against the enumeration of the boxes, multistart, CMA-ES, DIRECT and efficient
global optimization, and is applied to a feature-based topology optimization
problem with 108 design variables. The global optimum is reached at about one
third of the cost of the enumeration, and with fewer evaluations than the
baselines when the subdivision separates the basins of attraction. For the
topology optimization problem, the compliance is reduced from 84.0 to 82.1 with
respect to the solution obtained from the reference initial design.
